A bank that lends in Chile does not lend in pesos alone. A large share of mortgages, commercial loans, securities and bonds are written in the Unidad de Fomento (UF), a unit of account that the Banco Central de Chile calculates and publishes every day from the change of the consumer price index of the previous calendar month \citep{BCCh_UF}. A balance in UF is a balance in pesos that grows with measured inflation, so its repricing risk has two parts that a single curve cannot hold: the real rate at which the indexed flow is discounted and the indexation of its principal. Shiller describes the UF as the first successful indexed unit of account, in use in Chile since 1967 and later copied in Colombia, Ecuador, Mexico and Uruguay \citep{Shiller1998}. The supervisor of the Chilean banks, the Comisi\'on para el Mercado Financiero, therefore treats the flows indexed in local currency as a single currency of their own in the standard measures of interest rate risk in the banking book, called the UR, in Annex 1 of chapter 21-13 of its compilation of rules \citep{CMF_RAN_21_13}.

The Basel standard that the Chilean text adapts prescribes six rate scenarios applied currency by currency, with the size of each shock taken from a table of 21 currencies that holds neither the peso nor a unit of account indexed to prices \citep{BCBS_d368,BCBS_SRP31,BCBS_d578}. It leaves two things to national discretion that Annex 1 settles for the UF: a table of shock sizes for the domestic currency, and a floor on post-shock rates that must not exceed zero. Annex 1 gives the peso the sizes of 300, 400 and 200 basis points for the parallel, short and long shocks and gives the UR 200, 250 and 150 for the change in economic value of equity and 300 for the change in net interest income; it floors the post-shock peso rates at zero and the UR rates at minus two percent; it adds to the net interest income measure a charge of 2.5 percent of the absolute net UF position; and it combines the loss of one leg with the gain of the other by a pair rule that offsets half of the gain. A bank whose measures exceed 5 percent of level 1 capital for the income measure, or 15 percent for the value measure, or whose income measure exceeds 18 percent of its twelve-month margin, is examined in detail as a prioritised bank.

The question of this paper is what that separate treatment adds. A calculation that converted the UF flows to pesos at the UF value of the day, discounted them on the peso curve and shocked them with the peso table would be the natural single-currency reading of the Basel standard, and it is what a model that has no notion of an indexed unit of account does. The five components of the Chilean treatment (the curve, the shock sizes, the floor, the pair rule and the charge) act together in a measure, and a reader who sees only the sum of their effects cannot say which component carries a supervisory number. The supervisor's text does not say it, the literature on interest rate risk in banks is concerned with the exposure of real balance sheets, with deposit franchises and with hedging rather than with the construction of a standardised measure for an indexed unit \citep{Hoffmann2019,Drechsler2021,English2018}, and the Federal Reserve's review of the failure of Silicon Valley Bank, a failure of interest rate risk management, shows what is at stake when a measure understates an exposure \citep{FRB_SVB}. No public implementation of Annex 1 was found in the search behind this work, and it was not verified beyond that search: the claim of novelty is regional, and it is a claim about an open implementation that can be tested, not about the regulation, which is the supervisor's.

We propose a measurement. The treatment of the UF as its own currency is written as five explicit switches over one engine, so that the Annex 1 measure is the case with all five on and the single-currency treatment is the case with all five off. The 32 intermediate treatments are all computed, and the difference between the two measures is attributed to the switches by the exact Shapley value \citep{Shapley1953}, the decomposition that Shorrocks proposed for the attribution of a difference to its components \citep{Shorrocks2013}. A bound derived in the paper places the pair rule between the sum of the legs and the positive parts, a second result shows that the charge equals the worst loss of a symmetric surprise of 2.5 percentage points in inflation on the net UF position, and a third shows how a peso curve misvalues an indexed flow. The method runs on a synthetic bank, Banco Andino, calibrated to the aggregates of the banking system, over the real zero curves that the engine fits to the yields that the Banco Central de Chile publishes for the day \citep{BCCh_BDE}.

The contributions are the following.
\begin{enumerate}
\item An open implementation of the separate-currency treatment as five switches, run on the engine \texttt{riskvalidation} 1.1.0, whose Annex 1 measures reproduce the committed delta NII, delta EVE, charge and twelve-month margin of the six variants of the case to a relative deviation of \repMaxDev{} (Section~\ref{sec:results-reproduction}).
\item A bound, proved in Section~\ref{sec:theory}, that sandwiches the pair rule between two aggregation rules, and a counterexample in the data to the statement that the Annex 1 income measure always exceeds the Basel positive-part aggregation.
\item The exact Shapley decomposition of the gap between the two treatments into the five components, with 95 percent intervals over 30 synthetic books (Section~\ref{sec:ablation}).
\item A measurement of when the two treatments classify a bank differently under the three prioritisation tests, over net UF positions in the range from \sweepLo{} to \sweepHi{} percent of level 1 capital (Section~\ref{sec:results-sweep}).
\item A demonstration of the effect of the UR floor, which binds on no committed curve, by translating the real curve down, with the result that the floor can raise a measure (Section~\ref{sec:results-floors}).
\end{enumerate}

The main result is quantitative and limited. On the published case (seed 42) the single-currency delta NII is \repRatioLo{} to \repRatioHi{} times smaller than the Annex 1 one, and the single-currency delta EVE is \repEveShortLo{} to \repEveShortHi{} percent smaller, in the three distinct variants with a net UF position of 30 percent of level 1 capital; over 30 books the ratio of the mean income measures is $\rmRatioLng$, $\rmRatioSht$ and $\rmRatioOff$ (bootstrap 95 percent intervals $\rmLoLng$ to $\rmHiLng$, $\rmLoSht$ to $\rmHiSht$ and $\rmLoOff$ to $\rmHiOff$) and the single-currency delta EVE is $\mainEveShortLng\pm\mainEveShortLngCi$, $\mainEveShortSht\pm\mainEveShortShtCi$ and $\mainEveShortOff\pm\mainEveShortOffCi$ percent smaller. The indexation charge supplies about half of the Annex 1 income measure in those variants, while for the value measure the own curve and the own shock sizes are the large components, with opposite signs, and the pair rule is small except in the offset variant, where it carries \pairShareEveOff{} percent of the gap. The bank is synthetic, so the numbers measure the treatment on a plausible balance sheet and not the exposure of an existing bank. Section~\ref{sec:related-work} places the work among the published methods, Section~\ref{sec:problem} states the problem and Section~\ref{sec:proposed-method} the method, Section~\ref{sec:theory} proves the bounds, Section~\ref{sec:experimental-design} states the experiments, Sections~\ref{sec:results} and~\ref{sec:ablation} give the results and the ablation, and Sections~\ref{sec:discussion} to~\ref{sec:conclusions} close.
